% concatenated from A_labeling_of_the_Simplex-Lattice_Hypergraph_with_at_most_2_colors_on_each_hyperedge.tex
\documentclass[12pt]{article}

\usepackage{amssymb}
\usepackage{amsmath}
\usepackage{amsthm}
\usepackage{bm}



\newtheorem{theorem}{Theorem}
\newtheorem{proposition}{Proposition}
\newtheorem{definition}{Definition}
\newtheorem{lemma}{Lemma}
\newtheorem{corollary}{Corollary}




\begin{document}

\title{\Large A labeling of the Simplex-Lattice Hypergraph with at most 2 colors on each hyperedge}
\author{Ognjen Papaz\thanks{Faculty of Philosophy, University of East Sarajevo, Bosnia and Herzegovina {\tt ognjen.papaz@ff.ues.rs.ba}}
 \and Du\v sko Joji\'c\thanks{Faculty of Science, University of Banja Luka, Bosnia and Herzegovina {\tt dusko.jojic@pmf.unibl.org}}}
\date{}
\maketitle

\begin{abstract}
This paper provides a positive answer to the question of Mirzakhani and Vondrak from \cite{MV} that asks if there is a Sperner-admissible labeling of the simplex-lattice hypergraph such that each hyperedge uses at most 2 colors.
\end{abstract}

\section{Introduction}

In (\cite{MV}, section 4) the authors proved that for $k\geq 4$ and $q\geq k^2$, there is a Sperner-admissible labeling of the simplex-lattice hypergraph $H_{k,q}$ such that every hyperedge of $H_{k,q}$ contains at most 4 colors.

They raised the question of whether there is a Sperner-admissible labeling  of $H_{k,q}$ such that every hyperedge of $H_{k,q}$ contains at most 2 colors, motivation being provided by the fact that an answer would have consequences for fair division, see (\cite{MV}, sections 7 and 8).  



\section{Preliminaries}

We review basic definitions following \cite{MV}. 

Let $k\geq 3$ and $q\geq 1$ throughout the paper.

Denote with $R_{k,q}$ the simplex in $\mathbb R^{k-1}$ whose vertices are
\[(0,0,\ldots,0),(0,0,\ldots,0,q),\ldots,(0,q,q\ldots,q),(q,q,\ldots,q).\]

Let $V_{k,q}$ be the set of integer points in $R_{k,q}$, i.e.
\[V_{k,q}=\{\bm{v}=(v_1,v_2,\ldots,v_{k-1})\in \mathbb Z^{k-1}:0\leq v_i\leq v_{i+1}\leq q\}.\]

\textbf{The simplex-lattice hypergraph} is a $k$-uniform hypergraph $H_{k,q}$ whose vertex set is $V_{k,q}$ and whose set of hyperedges $E_{k,q}$ is given with
\[E_{k,q}=\{F(\bm{v}):\bm{v}\in V_{k,q-1}\},\]
where 
\[F(\bm{v})=\{\bm{v},\bm{v}+e_{k-1},\bm{v}+e_{k-1}+e_{k-2},\ldots,\bm{v}+e_{k-1}+e_{k-2}+\cdots+e_1\}.\]

The hyperedges of the simplex-lattice hypergraph $H_{k,q}$ correspond with the certain cells of the edgewise subdivision of the simplex $R_{k,q}$.\\

\textbf{The edgewise subdivision} of $R_{k,q}$ is its triangulation whose vertex set is $V_{k,q}$. The facets (maximal cells) of the triangulation are indexed with pairs $(\bm{v},\pi)$ where $\bm{v}\in V_{k,q-1}$ and $\pi\in\mathbb S_{k-1}$ is \textit{consistent} with $\bm{v}$. A permutation $\pi\in\mathbb S_{k-1}$ is consistent with $\bm{v}$ if $i$ appears before $i+1$ in $\pi$ whenever $v_i=v_{i+1}$. For each $\bm{v}\in V_{k,q-1}$ and each $\pi\in\mathbb S_{k-1}$ which is consistent with $\bm{v}$ the convex hull of the set
\[F(\bm{v},\pi)=\{\bm{v},\bm{v}+e_{\pi(k-1)},\bm{v}+e_{\pi(k-1)}+e_{\pi(k-2)},\ldots,\bm{v}+e_{\pi(k-1)}+e_{\pi(k-2)}+\cdots+e_{\pi(1)}\}\]
is a facet of the triangulation.

More details about the edgewise subdivision of a simplex can be found in \cite{ED}.\\

We see that each hyperedge $F(\bm{v})$ in $E_{k,q}$ is equal to $F(\bm{v},Id(\in\mathbb S_{k-1}))$. Extending this correspondence we define the $\pi$-simplex-lattice hypergraph.\\

\textbf{The $\pi$-simplex lattice hypergraph} is a $k$-uniform hypergraph $H_{k,q}^\pi \ (\pi\in\mathbb S_{k-1})$ whose set of hyperedges $E_{k,q}^\pi$ is given with
\[E_{k,q}^\pi=\{F(\bm{v},\pi):\bm{v}\in V_{k,q-1}, \pi\ \text{is consistent with}\ \bm{v}\}.\]

\textbf{The Sperner-admissible labeling} of the vertices $V_{k,q}$ is a mapping $\ell :V_{k,q}\to[k]$ such that $v_i>v_{i-1}$ whenever $\ell(\bm{v})=i$. 

Here we introduce the convention that $v_0=0$ and $v_k=q$.\\

 




\section{The labeling} 

\begin{definition}For each $\bm{v}\in V_{k,q}$ we define $r(\bm{v})$ and $i(\bm{v})$ in the following way:
\[r(\bm{v})=\max\{t-v_t:t\in[0,k]\},\]
\[i(\bm{v})=\min\{t\in[0,k]:t-v_t=r(v)\}.\]
Recall that $v_0=0$ and $v_k=q$.
\end{definition}

We can immediately see that $r(\bm{v})\geq v_0-0=0$ and that $i(\bm{v})=0$ if $r(\bm{v})=0$.\\

\begin{definition} We consider the mapping $\ell:V_{k,q}\to[k]$ defined with
\[\ell(\bm{v})=i(\bm{v})+1.\]
\end{definition}

\begin{proposition}
For $q>k$, the mapping $\ell$ is well-defined Sperner-admissible labeling of $V_{k,q}$.
\end{proposition}

\begin{proof} Let $\bm{v}\in V_{k,q}$ and $i=i(\bm{v})$. We first prove that $l(\bm{v})=i+1\in[k]$. Since $q>k$ we have $k-v_k=k-q<0$, hence $i<k$. Now we prove that $v_i<v_{i+1}$. By the definition of $i(\bm{v})$ we have that
\[i-v_i\geq i+1-v_{i+1},\]
hence $v_{i+1}\geq v_i+1$.
\end{proof}

\begin{theorem}
For $q>k$ and for the Sperner-admissible labeling $\ell$ each hyperedge of $H_{k,q}$ uses at most 2 colors. 
\end{theorem}

\begin{proof}
Let $F(\bm{v})$ be a hyperedge of $H_{k,q}$ and let $i=i(\bm{v})$ and $r=r(\bm{v})$. We denote the vertices of $F(\bm{v})$ in the following way
\[\bm{v}=\bm{v}, \bm{v}^{(k-1)}=\bm{v}+e_{k-1},\bm{v}^{(k-2)}=\bm{v}+e_{k-1}+e_{k-2},\ldots,\bm{v}^{(1)}=\bm{v}+e_{k-1}+e_{k-2}+\cdots+e_1.\]
Note that the vertices $\bm{v}^{(k-1)}, \bm{v}^{(k-2)},\ldots,\bm{v}^{(1)}$ are obtained by increasing coordinates of $\bm{v}$, one by one, from behind.

Let $r=0$, then $i=0$.  Since increasing the coordinates of $\bm{v}$ doesn't increase $r(\bm{v})$ we have that
\[0=r(\bm{v}^{(k-1)})=r(\bm{v}^{(k-2)})=\ldots=r(\bm{v}^{(1)}),\]
\[0=i(\bm{v}^{(k-1)})=i(\bm{v}^{(k-2)})=\ldots=i(\bm{v}^{(1)})\]
and
\[1=\ell(\bm{v})=\ell(\bm{v}^{(k-1)})=\ldots=\ell(\bm{v}^{(1)}).\]

Now suppose that $r>0$, then $i>0$. Increasing the coordinates $v_{i+1},\ldots,v_{k-1}$ of $\bm{v}$ doesn't change $r(\bm{v})$ or $i(\bm{v})$, hence
\[r=r(\bm{v}^{(k-1)})=r(\bm{v}^{(k-2)})=\ldots=r(\bm{v}^{(i+1)}),\]
\[i=i(\bm{v}^{(k-1)})=i(\bm{v}^{(k-2)})=\ldots=i(\bm{v}^{(i+1)}),\]
and
\[i+1=\ell(\bm{v})=\ell(\bm{v}^{(k-1)})=\ldots=\ell(\bm{v}^{(i+1)}).\]

Let's take a closer look at the vertex $\bm{v}^{(i)}$. By the definition of $r(\bm{v})$ we have
\[i-1-v^{(i)}_{i-1}=i-1-v_{i-1}\leq r(\bm{v}^{(i)})<r=i-v_i.\]
From here we see that $v_i<v_{i-1}+1$, hence $v_i=v_{i-1}$ and $r(\bm{v}^{(i)})=r-1$. 

Let $i'=i(\bm{v}^{(i)})$, this is the smallest index $t\in[0,k]$ such that $t-v_t=r-1$. We have 
\[r-1=r(\bm{v}^{(i)})=r(\bm{v}^{(i-1)})=\ldots=r(\bm{v}^{(i'+1)}),\]
\[i'=i(\bm{v}^{(i)})=i(\bm{v}^{(i-1)})=\ldots=i(\bm{v}^{(i'+1)}),\]
and  
\[i'+1=\ell(\bm{v}^{(i)})=\ell(\bm{v}^{(i-1)})=\ldots=\ell(\bm{v}^{(i'+1)}).\]

If $r-1=0$ then $i'=0$ and we are done. Let $r-1>0$, then $i'>0$. When we increase the coordinates $v_{i'},v_{i'+1},\ldots,v_{i},\ldots,v_{k-1}$ of $\bm{v}$ and obtain $\bm{v}^{(i')}$ we can see that $r(\bm{v}^{(i')})=r-1$ and that $i$ becomes the smallest index $t\in[0,k]$ such that $t-v^{(i')}_t=r-1$. Thus, 
\[r-1=r(\bm{v}^{(i')})=r(\bm{v}^{(i'-1)})=\ldots=r(\bm{v}^{(1)}),\]
\[i=i(\bm{v}^{(i')})=i(\bm{v}^{(i'-1)})=\ldots=i(\bm{v}^{(1)}),\]
and
\[i+1=\ell(\bm{v}^{(i')})=\ell(\bm{v}^{(i'-1)})=\ldots=\ell(\bm{v}^{(1)}).\]
\end{proof}

Now we consider if there exist a Sperner-admissible labeling of $V_{k,q}$ such that each hyperedge of $H_{k,q}^\pi\ (\pi\in\mathbb S_{k-1})$ uses at most 2 colors.\\

Analyzing the proof of the previous theorem we see that the labeling $\ell$ works because among all coordinates $v_t$ of $\bm{v}$ such that $t-v_t=r(\bm{v})$ the coordinate $v_{i(\bm{v})}$ is the last coordinate that increases when obtaining the vertices of the hyperedge $F(\bm{v})$.

In a hyperedge 
\[F(\bm{v},\pi)=\{\bm{v},\bm{v}+e_{\pi(k-1)},\bm{v}+e_{\pi(k-1)}+e_{\pi(k-2)},\ldots,\bm{v}+e_{\pi(k-1)}+e_{\pi(k-2)}+\cdots+e_{\pi(1)}\}\]
of $H_{k,q}^\pi$, the permutation $\pi$ prescribes the order in which the coordinates of $\bm{v}$ are increasing.
 Thus, we can modify the definitions of $i(\bm{v})$ and $\ell$ accordingly and show that each hyperedge of $H_{k,q}^\pi$ will use at most two colors by the same arguments as for $\ell$.

\begin{definition} For $\pi\in\mathbb {S}_{k-1}$, let $\overline\pi=0\pi k$. For each $\bm{v}\in V_{k,q}$ we define $i^\pi(\bm{v})$ with 
\[i^\pi(\bm{v})=\min\{t\in[0,k]:\overline\pi(t)-v_{\overline\pi(t)}=r(\bm{v})\}\]
and $\ell^\pi(\bm{v})$ with $\ell^\pi(\bm{v})=i^\pi(\bm{v})+1$. 
\end{definition}

\begin{theorem}
For $q>k$ the mapping $\ell^\pi$ is a Sperner-admissible labeling such that each hyperedge of $H_{k,q}^\pi$ uses at most 2 colors.
\end{theorem}

\bibliographystyle{plain}

\begin{thebibliography}{10}

\bibitem{ED}

Herbert Edelsbrunner  and Daniel R. Grayson.
\newblock Edgewise subdivision of a simplex.
\newblock {\em Discrete and Computation Geometry}, 24(4):707-719,2000


\bibitem{MV}
Maryam Mirzakhani and Jan Vondr{\'{a}}k.
\newblock Sperner's colorings, hypergraph labeling problems and fair division.
\newblock {\em Proc. of ACM-SIAM SODA}, 873--886, 2015.



\end{thebibliography}




























\end{document}

